\documentclass[10pt,a4paper]{amsart}
\usepackage[margin=19mm]{geometry}
\usepackage{amsmath,amssymb,mathtools}
\usepackage[scr=rsfs,cal=euler]{mathalfa}
\usepackage{xfrac}
\usepackage[unicode,hidelinks,bookmarks=false]{hyperref}
\newcommand{\ie}{i.e.\ }
\newcommand{\cf}{cf.\ }
\newcommand{\resp}{respectively\ }
\newcommand{\Subsection}{\subsection}
\providecommand{\nobreakdash}{\nobreak\hbox{-}}
\setlength{\emergencystretch}{2em}
\multlinegap0pt
\usepackage{bm}
\usepackage{bbm}
\usepackage{dsfont}
\usepackage{xspace}
\usepackage{cancel}
\usepackage{mathtools}
\usepackage{cleveref}
\usepackage[shortlabels]{enumitem}
\usepackage{amsthm}
\makeatletter
\@ifundefined{theorem}{\newtheorem{theorem}{Theorem}}{}
\@ifundefined{lemma}{\newtheorem{lemma}[theorem]{Lemma}}{}
\makeatother
\title[Convergence of Pivoted Cholesky Algorithm for Lipschitz Kern]{Convergence of Pivoted Cholesky Algorithm for Lipschitz Kernels}
\author{Sungwoo Jeong, Alex Townsend}
\date{Source: arXiv:2509.13582v2 (CC BY 4.0)}
\begin{document}
\maketitle

\noindent\emph{Source and licence.} Convergence of Pivoted Cholesky Algorithm for Lipschitz Kernels. Version arXiv:2509.13582v2 (CC BY 4.0), distributed under the Creative Commons Attribution 4.0 International licence, as recorded in the arXiv licence metadata for this exact version and shown on the abstract page https://arxiv.org/abs/2509.13582v2. Full rights notice: licenses/arxiv-2509.13582-CC-BY-4.0.txt.\par\smallskip

\noindent\textit{Editorial scope.} Counted unit: the Lemma whose source label is \texttt{lem:lipschitzchangemulti} in the article (source0.tex lines 434--456, source label \texttt{lem:lipschitzchangemulti}), delivered together with the article's own complete original proof of it (source0.tex lines 457--487, one \texttt{proof} environment of 31 lines). No mathematics is written, repaired or re-derived: statement and proof are the article's own text line for line. Required context retained uncounted: none. Every definition, display and label of those lines is retained unchanged. Omitted from this package and not redistributed: 1--433, 488--720. Nothing inside a retained excerpt is omitted or altered: every retained body is delivered exactly as the article's lines are written, so no omission receipt of any kind accompanies this package. Citations: the retained excerpts carry no citation command at all, so no bibliography entry of the article is transcribed and no reference is redistributed. Reference adaptation: none was needed and none was made; every reference inside the delivered unit resolves inside it, so no unresolved reference or citation is produced. Printed slips kept literal: no printed word, symbol, hypothesis or number of the counted statement or of its proof was altered. Layout and class: the article's own class is \texttt{amsart} with packages including cleveref, graphicx, geometry, amssymb, amsmath, amsfonts, enumitem, comment, overpic, amsaddr, xcolor, todonotes; the wrapper is the lane's pinned \texttt{amsart} editorial wrapper at 10pt on A4, which restates the theorem-like environments the retained text needs and adds the article's own macro definitions that the retained text needs (none), with extra wrapper packages (bm, bbm, dsfont, xspace, cancel, mathtools, cleveref, enumitem). The delivered numbering is the unit's own. Assets: no retained span contains a figure environment, a graphics-inclusion command, a PDF-page inclusion, a table or a TikZ picture; no image file is copied into this package, the package declares no asset dependency and no image file is redistributed with it. Licence: version arXiv:2509.13582v2 of this article is distributed under the Creative Commons Attribution 4.0 International licence, as recorded in the arXiv licence metadata for this exact version and shown on the abstract page https://arxiv.org/abs/2509.13582v2. A full rights notice is delivered at licenses/arxiv-2509.13582-CC-BY-4.0.txt. Characters: every retained excerpt is pure ASCII, so the delivered engine, XeLaTeX with its default Latin Modern text font, reports no missing character.
\begin{lemma}\label{lem:lipschitzchangemulti}
Let $\Omega\subset\mathbb{R}^d$ be compact and let $K:\Omega\times\Omega\to\mathbb{R}$ be a continuous SPD kernel.
Assume:
\begin{enumerate}[label=(\alph*)]
\item There exists $L>0$ such that, for all $1\le j\le d$,
\begin{equation}\label{eq:diag-deriv-lip}
\bigl|\partial_j^{1}K(x,y)-\partial_j^{1}K(x,x)\bigr|\le L\|x-y\|,
\qquad
\bigl|\partial_j^{2}K(x,y)-\partial_j^{2}K(x,x)\bigr|\le L\|x-y\|,
\quad x,y\in\Omega,
\end{equation}
where the partials exist on $\Omega\times\Omega$.
\item There exists $L_0>0$ such that
\begin{equation}\label{eq:first-arg-lip}
\bigl|K(x,y)-K(x',y)\bigr|\le L_0\|x-x'\|\qquad \text{for all }x,x',y\in\Omega,
\end{equation}
and $K$ is differentiable in its first argument on $\Omega\times\Omega$.
\end{enumerate}
Let $R_1$ be the residual after one Cholesky step (with any pivoting strategy) and let $z_1$ be the first pivot. Then the diagonal of the residual $D_1(x)=R_1(x,x)$ is differentiable on $\Omega$, and for each $j=1,\dots,d$ there exists a constant $C>0$ (depending only on $L$, $L_0$, and $K_{\min}=\min_{x\in\Omega}K(x,x)>0$) such that
\begin{equation}\label{eq:grad-D1-linear}
\bigl|\partial_j D_1(x)\bigr|\;\le\; \left(4L + \frac{2L_0^2}{K(z_1,z_1)}\right)\,\|x-z_1\|\qquad \text{for all }x\in\Omega.
\end{equation}
\end{lemma}

\begin{proof}
Write
\[
D_1(x)=K(x,x)-\frac{K(x,z_1)^2}{K(z_1,z_1)},\qquad
\partial_j D_1(x)
= \partial_j^1K(x,x)+\partial_j^2K(x,x)
-\frac{2}{K(z_1,z_1)}\,K(x,z_1)\,\partial_j^1K(x,z_1).
\]
Because $R_1(\cdot,z_1)=R_1(z_1,\cdot)=0$, we have $\partial_j D_1(z_1)=0$ for all $j$.
Subtracting $\partial_j D_1(z_1)$ and using the triangle inequality,
\begin{align*}
|\partial_j D_1(x)|
&= \bigl|\partial_j D_1(x)-\partial_j D_1(z_1)\bigr| \\
&\le \underbrace{\bigl|\partial_j^1K(x,x)-\partial_j^1K(x,z_1)\bigr|}_{\le\, L\|x-z_1\|\ \text{by }\eqref{eq:diag-deriv-lip}}
+\underbrace{\bigl|\partial_j^2K(x,x)-\partial_j^2K(x,z_1)\bigr|}_{\le\, L\|x-z_1\|\ \text{by }\eqref{eq:diag-deriv-lip}} \\
&\quad + \frac{2}{K(z_1,z_1)}
\underbrace{\bigl|K(x,z_1)-K(z_1,z_1)\bigr|}_{\le\, L_0\|x-z_1\|\ \text{by }\eqref{eq:first-arg-lip}}
\ \underbrace{\bigl|\partial_j^1K(x,z_1)\bigr|}_{\le\, L_0\ \text{since }K\text{ is }L_0\text{-Lipschitz in }x} \\
&\quad + \frac{2}{K(z_1,z_1)}\,K(z_1,z_1)\,
\underbrace{\bigl|\partial_j^1K(x,z_1)-\partial_j^1K(z_1,z_1)\bigr|}_{\le\, L\|x-z_1\|\ \text{by }\eqref{eq:diag-deriv-lip}}.
\end{align*}
Here we used that \cref{eq:first-arg-lip} and differentiability in $x$ imply
\(
\sup_{(x,y)\in\Omega^2}|\partial_j^{1}K(x,y)|\le L_0
\)
(by the mean value theorem). Combining the bounds gives
\[
|\partial_j D_1(x)| \;\le\; \left(4L + \frac{2L_0^2}{K(z_1,z_1)}\right)\,\|x-z_1\|,
\]
which proves \cref{eq:grad-D1-linear}.
\end{proof}

\end{document}
